\chapter{Zaklju\v{c}ak}
\label{ch:zakljucak}


\section{Rezime (osvrt na ciljeve iz zadatka)}

The goal set in the assignment was to implement four ray-traced lighting effects, shadows, ambient
occlusion, reflections, and diffuse global illumination, through inline ray queries in a custom
Vulkan engine, and to characterize the cost and quality of each. All four effects were implemented as
described in Chapter~\ref{ch:implementacija}, sharing one hit-shading routine over a bindless geometry
table and one reusable SVGF denoiser, and integrated into an existing rasterizing forward renderer
rather than a standalone tracer. The evaluation in Chapter~\ref{ch:analiza} isolates the per-effect
GPU cost with an A/B benchmark and measures reconstruction quality against a converged path-traced
reference, and compares the identical pipeline
across two GPU vendors (AMD RDNA4 and NVIDIA Blackwell).
The headline numbers are collected in Table~\ref{tab:summary}. The two cards do not agree on which
effect costs most: on the RX 9070 XT shadows dominate, while on the RTX 5070 shadows, reflections and
global illumination sit within \perfNVSpread{} ms of one another, which is inside the run-to-run
spread, so on that card they are indistinguishable in cost.

\begin{table}[tbp]
\caption[Sa\v{z}etak izmerenih rezultata]{Sa\v{z}etak izmerenih rezultata. Cena po efektu je razlika
ukupnog GPU vremena izme\dj u susednih konfiguracija; kvalitet je prosek preko tri pogleda naspram
konvergirane path-traced reference na RX 9070 XT.}
\label{tab:summary}
\centering
\begin{tabular}{lrr}
\toprule
 & RX 9070 XT & RTX 5070 \\
\midrule
osnova bez RT (ms)      & \perfBaseAMD{}        & \perfBaseNV{} \\
senke (ms)              & \perfEffShadowsAMD{}  & \perfEffShadowsNV{} \\
AO (ms)                 & \perfEffAoAMD{}       & \perfEffAoNV{} \\
refleksije (ms)         & \perfEffReflAMD{}     & \perfEffReflNV{} \\
GI (ms)                 & \perfEffGiAMD{}       & \perfEffGiNV{} \\
\midrule
ukupno, sva \v{c}etiri (ms) & \perfTotalAMD{}   & \perfTotalNV{} \\
\midrule
FLIP: raster $\rightarrow$ sve RT  & \multicolumn{2}{c}{\qRasterFlip{} $\rightarrow$ \qAllRtFlip{}} \\
PSNR (dB): raster $\rightarrow$ sve RT & \multicolumn{2}{c}{\qRasterPsnr{} $\rightarrow$ \qAllRtPsnr{}} \\
SSIM: raster $\rightarrow$ sve RT  & \multicolumn{2}{c}{\qRasterSsim{} $\rightarrow$ \qAllRtSsim{}} \\
\bottomrule
\end{tabular}
\end{table}

Two of the sweeps sharpen that picture. The ray-count knee sits at one ray per pixel: going from one
to eight rays costs \rcGpuPct\% more GPU time for a \rcPsnrGain{} dB PSNR gain, because the denoiser and temporal
accumulation already recover most of the achievable quality from a single sample
(Section~\ref{sec:raycount}). The spatial denoiser's own contribution is likewise not where a static
metric can see it: it moves PSNR, SSIM and FLIP by less than run-to-run noise, while cutting
frame-to-frame RMS on a static view from \flickerOff{} to \flickerOn{}, a \flickerPct\% reduction, for \perfDenoiseAMD{} ms
(Section~\ref{sec:denoise-eval}). The honest reading is that at this sample count the temporal history
does the reconstruction and the spatial filter buys stability rather than static fidelity.

The practical value is a demonstration that a hybrid
ray-query pipeline gives correct off-screen lighting at an interactive budget without a full path
tracer, which is directly applicable to real-time engines that want physically grounded indirect
lighting, contact shadows, and reflections.

\section{Ograni\v{c}enja}

The work has clear limitations, stating them sharpens what the results do and do not show. The
measurements are taken on two GPUs of comparable market tier and a single scene class, so the absolute numbers do not
generalize across hardware or content. Global illumination is limited to a single diffuse bounce, and
shadows are hard or lightly softened rather than physically accurate area shadows. To stay within a
real-time budget the ray counts are kept very low (one or two per pixel), which places heavy demand
on the denoiser and exposes its failure modes: temporal ghosting and lag under fast motion,
particularly for reflections, which reuse the diffuse-tuned temporal reprojection. The
top-level acceleration structure is fully rebuilt on every dynamic frame rather than refit, which is
heavier than necessary for small transform changes. Finally, using inline ray queries
gives up the hardware ray scheduling and recursion that a ray-tracing pipeline would provide, a
deliberate tradeoff for integration simplicity that would matter for deeper or more numerous bounces.

\section{Budu\'ci rad}

Several directions extend the work. Sampling is the most immediate: replacing the naive per-light and
cosine-hemisphere sampling with reservoir resampling (ReSTIR and ReSTIR GI) would improve quality at
equal ray count, especially for many lights and for indirect paths
\cite{bitterli2020restir, ouyang2021restirgi}. Reconstruction is the second: the engine already
contains a temporal neural network used as an anti-aliasing resolve, and a natural next step is a
learned denoiser or upscaler replacing or augmenting SVGF, together with the vendor-neutral,
cross-vendor cooperative-matrix inference study that motivated an earlier version of this project. A
third direction is scaling the denoised shadows to many lights more cheaply, with a MegaLights-style
stochastic per-pixel light selection on the shared denoiser, rather than tracing every light every
frame; a dedicated, penumbra-aware shadow denoiser in place of the shared one is a related refinement. Finally, the cross-vendor comparison here covers two current-generation cards; broadening it to more cards and
older generations, plus a second rendering backend (DirectX 12), would extend the performance story further.
